\documentclass[12pt, a4paper]{article}

\usepackage{amsmath, amssymb}

\usepackage{fontspec}
\setmainfont{FreeSerif} % FreeSerif-ը Overleaf-ում լռելյայն ունի հայերենի գերազանց աջակցություն

\usepackage{geometry}
\geometry{margin=1in}

\usepackage[colorlinks=true, linkcolor=blue, urlcolor=blue, citecolor=blue]{hyperref}

\setcounter{secnumdepth}{0} 

\begin{document}

\tableofcontents
\vspace{2em}
\hrule
\vspace{2em}

\section{1. Անկյունագծայնացման տեսություն և հատկություններ}

\subsection{1.1. Սահմանում}

$A \in M_{n \times n}$ քառակուսի մատրիցը կոչվում է \textbf{անկյունագծայնացվող}, եթե այն \textbf{նման է} որևէ $D$ անկյունագծային մատրիցի։ Ըստ սահմանման, գոյություն ունի այնպիսի հակադարձելի $P$ մատրից, որ.

$$
A = P D P^{-1}
$$

որտեղ $D = \text{diag}(\lambda_1, \lambda_2, \dots, \lambda_n)$ մատրիցի գլխավոր անկյունագծի տարրերը $A$-ի սեփական արժեքներն են, իսկ $P$ մատրիցի սյուները՝ դրանց համապատասխանող գծորեն անկախ սեփական վեկտորները։

\subsection{1.2. Նման մատրիցների հատկությունները}

Եթե $A \sim B$, ապա նրանք օժտված են նույն սեփական արժեքներով, քանի որ նման մատրիցների բնութագրիչ բազմանդամները նույնական են.

$$
\det(A - \lambda I) = \det(B - \lambda I)
$$

\subsection{1.3. Հակադարձ մատրիցի սեփական արժեքները}

\textbf{Պնդում.} Եթե $\lambda \neq 0$-ն $A$ մատրիցի սեփական արժեքն է, ապա $\lambda^{-1}$-ը $A^{-1}$ մատրիցի սեփական արժեքն է։

\textbf{Ապացույց.}\\
Դիցուք $Ax = \lambda x$, որտեղ $x \neq 0$։ Քանի որ $A$-ն հակադարձելի է ($\det A \neq 0 \implies \lambda \neq 0$), հավասարման ձախակողմյան բազմապատկումը $A^{-1}$-ով տալիս է.

$$
A^{-1}(Ax) = A^{-1}(\lambda x) \implies x = \lambda (A^{-1}x) \implies A^{-1}x = \frac{1}{\lambda}x
$$

Այսպիսով, $x$-ը հանդիսանում է $A^{-1}$-ի սեփական վեկտոր $\frac{1}{\lambda}$ սեփական արժեքով։

\vspace{1em}
\hrule
\vspace{1em}

\section{2. Գործնական վարժություններ և լուծումներ}

\subsection{Խնդիր 253. Անկյունագծային տեսքի բերում}

\textbf{Օրինակ ա):}

$$
A = \begin{pmatrix}
-14 & 12 \\
-20 & 17
\end{pmatrix}
$$

\begin{enumerate}
    \item \textbf{Բնութագրիչ հավասարում.} $\lambda^2 - 3\lambda + 2 = 0 \implies \lambda_1 = 1, \lambda_2 = 2$։
    \item \textbf{Սեփական վեկտորներ.}
\end{enumerate}

$$
v_1 =
\begin{pmatrix}
4 \\
5
\end{pmatrix},
\quad
v_2 =
\begin{pmatrix}
3 \\
4
\end{pmatrix}
$$

\begin{enumerate}
    \setcounter{enumi}{2}
    \item \textbf{Մատրիցների կազմում.}
\end{enumerate}

$$
P =
\begin{pmatrix}
4 & 3 \\
5 & 4
\end{pmatrix},
\quad
P^{-1} =
\begin{pmatrix}
4 & -3 \\
-5 & 4
\end{pmatrix}
$$

\begin{enumerate}
    \setcounter{enumi}{3}
    \item \textbf{Արդյունք.}
\end{enumerate}

$$
P^{-1} A P =
\begin{pmatrix}
1 & 0 \\
0 & 2
\end{pmatrix}
$$

\begin{quote}
\textbf{Հակադարձ մատրիցի հաշվարկման մեթոդներ.}

\begin{enumerate}
    \item \textbf{$2 \times 2$ դեպքում.}
\end{enumerate}

$$
A^{-1} = \frac{1}{\det(A)}
\begin{pmatrix}
d & -b \\
-c & a
\end{pmatrix}
$$

\begin{enumerate}
    \setcounter{enumi}{1}
    \item \textbf{$n \times n$ դեպքում (Գաուս-Ժորդան).}
\end{enumerate}

$$
[P \mid I] \xrightarrow{\text{տողային ձևափոխություններ}} [I \mid P^{-1}]
$$
\end{quote}

\vspace{1em}
\hrule
\vspace{1em}

\section{3. Մատրիցի աստիճանի բարձրացման մեթոդաբանությունը}

\textbf{Խնդիր 256.} Հաշվել $A^{10}$, որտեղ

$$
A =
\begin{pmatrix}
1 & 0 \\
-1 & 2
\end{pmatrix}
$$

\begin{enumerate}
    \item \textbf{Սեփական արժեքներ.} Քանի որ $A$-ն ստորին եռանկյուն մատրից է՝ $\lambda_1 = 1, \lambda_2 = 2$։
    \item \textbf{Սեփական վեկտորներ.}
\end{enumerate}

$$
v_1 =
\begin{pmatrix}
1 \\
1
\end{pmatrix},
\quad
v_2 =
\begin{pmatrix}
0 \\
1
\end{pmatrix}
$$

$$
P =
\begin{pmatrix}
1 & 0 \\
1 & 1
\end{pmatrix}
$$

\begin{enumerate}
    \setcounter{enumi}{2}
    \item \textbf{Աստիճանի բարձրացում.} Օգտագործվում է $A^n = P D^n P^{-1}$ բանաձևը.
\end{enumerate}

$$
D^{10} =
\begin{pmatrix}
1^{10} & 0 \\
0 & 2^{10}
\end{pmatrix}
=
\begin{pmatrix}
1 & 0 \\
0 & 1024
\end{pmatrix}
$$

$$
A^{10} =
\begin{pmatrix}
1 & 0 \\
1 & 1
\end{pmatrix}
\begin{pmatrix}
1 & 0 \\
0 & 1024
\end{pmatrix}
\begin{pmatrix}
1 & 0 \\
-1 & 1
\end{pmatrix}
=
\begin{pmatrix}
1 & 0 \\
-1023 & 1024
\end{pmatrix}
$$

\vspace{1em}
\hrule
\vspace{1em}

\section{4. Սիմետրիկ մատրիցներ և SVD թեորեմը}

\subsection{4.1. Սպեկտրալ թեորեմ սիմետրիկ մատրիցների համար}

Եթե $A = A^T$, ապա գոյություն ունի $P$ օրթոգոնալ մատրից ($P^{-1} = P^T$) այնպես, որ.

$$
A = P D P^T
$$

\subsection{4.2. Եզակի արժեքների վերլուծություն (SVD)}

SVD-ն հանդիսանում է անկյունագծայնացման ընդհանրացումը կամայական $m \times n$ չափի մատրիցների համար։

\textbf{Թեորեմ (SVD).} Ցանկացած $A \in M_{m \times n}$ մատրիցի համար գոյություն ունեն $U \in M_{m \times m}$ և $V \in M_{n \times n}$ օրթոգոնալ մատրիցներ և $\Sigma \in M_{m \times n}$ ուղղանկյուն անկյունագծային մատրից այնպես, որ.

$$
A = U \Sigma V^T
$$

որտեղ $\Sigma_{ii} = \sigma_i \ge 0$ տարրերը կոչվում են $A$ մատրիցի եզակի արժեքներ։

\vspace{1em}
\hrule
\vspace{1em}

\section{5. SVD բաղադրիչների հաշվարկման տեսական հիմքերը}

$U$ և $V$ մատրիցները ստացվում են համապատասխանաբար $A A^T$ և $A^T A$ սիմետրիկ մատրիցների սպեկտրալ վերլուծությունից.

\begin{enumerate}
    \item \textbf{Աջ եզակի վեկտորներ ($V$).}
\end{enumerate}

$$
A^T A = (V \Sigma^T U^T)(U \Sigma V^T) = V (\Sigma^T \Sigma) V^T
$$

$V$-ի սյուները $A^T A$-ի սեփական վեկտորներն են։

\begin{enumerate}
    \setcounter{enumi}{1}
    \item \textbf{Ձախ եզակի վեկտորներ ($U$).}
\end{enumerate}

$$
A A^T = (U \Sigma V^T)(V \Sigma^T U^T) = U (\Sigma \Sigma^T) U^T
$$

$U$-ի սյուները $A A^T$-ի սեփական վեկտորներն են։

\begin{enumerate}
    \setcounter{enumi}{2}
    \item \textbf{Եզակի արժեքներ ($\sigma_i$).}
\end{enumerate}

$$
\sigma_i = \sqrt{\lambda_i(A^T A)}
$$

\vspace{1em}
\hrule
\vspace{1em}

\section{6. Մատրիցի ռանգի տեսություն}

\subsection{6.1. Ռանգի հիմնական հատկությունները}

\begin{itemize}
    \item $\text{rank}(A) = \text{rank}(A^T)$
    \item $\text{rank}(AB) \le \min(\text{rank}(A), \text{rank}(B))$
\end{itemize}

\subsection{6.2. Թեորեմ (Ռանգի մասին հիմնական թեորեմը)}

Ցանկացած $A \in M_{m \times n}$ մատրիցի համար սյունային ռանգը հավասար է տողային ռանգին։

\textbf{Ապացույցի սխեմա.}

\begin{enumerate}
    \item Եթե $\text{rank}_{col}(A) = r$, ապա $A = BC$, որտեղ $B \in M_{m \times r}$ և $C \in M_{r \times n}$։
    \item $A$-ի յուրաքանչյուր տող $C$-ի $r$ հատ տողերի գծային կոմբինացիա է $\implies \text{rank}_{row}(A) \le r$։
    \item Կիրառելով նույնը $A^T$-ի համար՝ ստանում ենք հակադարձ անհավասարությունը, ինչից բխում է հավասարությունը։
\end{enumerate}

\vspace{1em}
\hrule
\vspace{1em}

\section{7. SVD հաշվարկման ալգորիթմները}

\subsection{Եղանակ 1. Առանձնացված հաշվարկ}

$U$ և $V$ մատրիցները հաշվվում են անկախաբար՝ որպես $A A^T$ և $A^T A$ մատրիցների սեփական վեկտորների համախմբեր։

\subsection{Եղանակ 2. Կապակցված հաշվարկ}

\begin{enumerate}
    \item Հաշվել $V$ մատրիցը և $\sigma_i$ արժեքները։
    \item Զրոյից տարբեր $\sigma_i$-ների համար հաշվել
\end{enumerate}

$$
u_i = \frac{1}{\sigma_i} A v_i
$$

\begin{enumerate}
    \setcounter{enumi}{2}
    \item Եթե $\text{rank}(A) < m$, ապա լրացնել $U$ մատրիցը մինչև օրթոնորմալ բազիս՝ օգտագործելով $\text{null}(A^T)$ տարածությունը։
\end{enumerate}

\vspace{1em}
\hrule
\vspace{1em}

\section{8. Երկրաչափական մեկնաբանություն}

SVD-ն ցանկացած գծային ձևափոխություն տրոհում է երեք հաջորդական քայլերի.

\begin{enumerate}
    \item \textbf{$V^T$ (Պտույտ).} Մուտքային բազիսի կողմնորոշում ըստ $A^T A$-ի սեփական առանցքների։
    \item \textbf{$\Sigma$ (Ձգում).} Տարածության ձգում կամ սեղմում $\sigma_i$ գործակիցներով (միավոր շրջանագիծը վերածվում է էլիպսի)։
    \item \textbf{$U$ (Պտույտ).} Ելքային էլիպսի վերջնական կողմնորոշում թիրախային տարածությունում։
\end{enumerate}

\vspace{1em}
\hrule
\vspace{1em}

\section{9. Հավելյալ նշումներ}

\begin{itemize}
    \item \textbf{Տվյալների սեղմում.}
\end{itemize}

$$
A \approx \sum_{i=1}^{k} \sigma_i u_i v_i^T
$$

(Low-rank approximation)

\begin{itemize}
    \item \textbf{Օրթոգոնալության կարևորությունը.} $U$ և $V$ մատրիցների օրթոգոնալությունը երաշխավորում է, որ ձևափոխության երկրաչափական աղավաղումները ներառված են բացառապես $\Sigma$ մատրիցում։

    \item \textbf{Կարգավորվածություն.} $\Sigma$ մատրիցում եզակի արժեքները միշտ դասավորվում են նվազման կարգով՝
\end{itemize}

$$
\sigma_1 \ge \sigma_2 \ge \dots \ge \sigma_r > 0
$$

\end{document}